% 张三慧《大学物理学》第4版第16章：恒定电流，独立解稿
% 源教材V2 PDF物理页118--120；图16.19--16.24的电源极性仍待原图与辅导书复核。
% 所有结果为待审候选，不宣称专家认证。
\section*{第16章\quad 恒定电流：独立解析}

\paragraph{习题 16.1}
束流强度为 \(I\)，周长 \(L\)，电子速率
\(v\simeq c\)，则电子绕环周期
\(T=L/v\)。
在任意截面每秒通过电子数
\(I/e\)，同时在环中的电子数
\[
 \boxed{N=\frac Ie\frac Lc
 =\frac{(8.0\times10^{-3})(240)}
 {(1.602176634\times10^{-19})(3.00\times10^8)}
 \approx4.00\times10^{10}}.
\]
该估算忽略电子速度与 \(c\) 的微小差异和束团不均匀性。

\paragraph{习题 16.2}
假设绝缘传送带表面电荷密度 \(\sigma\)，
孤立大薄面每侧场强
\(E=\sigma/(2\varepsilon_0)\)。
于是
\(\sigma=2\varepsilon_0E
 \approx2.13\times10^{-5}\,\mathrm{C/m^2}\)。
单位时间传输的面积为带宽 \(w\)
乘速度 \(u\)，电流
\[
 \boxed{I=\sigma wu
 =2\varepsilon_0E(0.30)(0.20)
 \approx1.28\times10^{-6}\,\mathrm A
 =1.28\times10^{-3}\,\mathrm{mA}}.
\]
这是假设一面有效载电且外场取两侧薄面公式，
不是有导体背板时的局部电场边界模型。

\paragraph{习题 16.3}
直径 \(d=1.0\,\mathrm{mm}\)
的银丝截面积
\(S=\pi(d/2)^2=7.85\times10^{-7}\,\mathrm{m^2}\)。
用估计参数
\(\rho_{\rm Ag}\approx10.5\times10^3\,\mathrm{kg/m^3}\)，
\(M_{\rm Ag}\approx0.1079\,\mathrm{kg/mol}\)，
每银原子给一导电电子，则
\[
 n=\frac{\rho_{\rm Ag}N_A}{M_{\rm Ag}}
 \approx5.86\times10^{28}\,\mathrm{m^{-3}}.
\]
因此
\[
 \boxed{v_d=\frac{I}{neS}
 =\frac{30}{(5.86\times10^{28})e
 (7.85\times10^{-7})}
 \approx4.07\times10^{-3}\,\mathrm{m/s}}.
\]
经典理想气体近似下热运动平均\emph{速率}
\[
 \boxed{\bar v=\sqrt{\frac{8k_BT}{\pi m_e}}
 \approx1.06\times10^5\,\mathrm{m/s}}\quad(T=293\,\mathrm K).
\]
两者相差约七个数量级。
实际金属电子是简并费米气体，
此处经典热速率只是按原题要求给的模型估算，
不可声称它是银电子真实费米速度。

\paragraph{习题 16.4}
铜棒长 \(L=2\,\mathrm m\)，
截面积
\(S=(0.020)(0.080)=1.60\times10^{-3}\,\mathrm{m^2}\)，
施加电压 \(U=0.050\,\mathrm V\)。
电阻与电流为
\[
 \boxed{R=\rho L/S=2.1875\times10^{-5}\,\Omega},
\quad
 \boxed{I=U/R=2.286\times10^3\,\mathrm A}.
\]
电流密度、电场与耗散功率分别为
\[
 \boxed{j=I/S=1.429\times10^6\,\mathrm{A/m^2}},
\quad
 \boxed{E=U/L=0.0250\,\mathrm{V/m}},
\]
\[
 \boxed{P=UI=114.3\,\mathrm W}.
\]
自由电子的\emph{漂移速率}大小
\[
 \boxed{|v_d|=\frac{j}{ne}
 =\frac{1.429\times10^6}
 {(8.5\times10^{28})(1.6022\times10^{-19})}
 \approx1.05\times10^{-4}\,\mathrm{m/s}}.
\]
电子实际漂移方向与传统电流方向相反。

\paragraph{习题 16.5}
水和铁管沿相同轴向长度受到相同电压，
两者电流与各自电导截面积成正比。
铁质横截面积
\[
 S_{\rm Fe}=\frac\pi4
 [(0.020)^2-(0.015)^2]
 \approx1.374\times10^{-4}\,\mathrm{m^2},
\]
水截面积
\(S_w=\pi(0.015)^2/4
 \approx1.767\times10^{-4}\,\mathrm{m^2}\)。
所以
\[
 \frac{I_w}{I_{\rm Fe}}=
 \frac{S_w/\rho_w}{S_{\rm Fe}/\rho_{\rm Fe}}
 =\frac{S_w\rho_{\rm Fe}}{S_{\rm Fe}\rho_w}
 \approx1.12\times10^{-5}.
\]
合计电流为20 A，故
\[
 \boxed{I_w\approx2.24\times10^{-4}\,\mathrm A},
\qquad
 \boxed{I_{\rm Fe}\approx19.9998\,\mathrm A}.
\]
水的电阻率按题给 \(0.01\,\Omega\mathrm m\)；
并不意味着所有实际自来水都有这一固定值。

\paragraph{习题 16.6}
设短路点至左端距离 \(x\)，
每根均匀线的单位长度电阻 \(r'\)，
从左端AB测量有两根导线从左端
通到故障点再返回，
\[
 R_{AB}=2r'x=30\,\Omega.
\]
从右端CD测量
\(R_{CD}=2r'(L-x)=70\,\Omega\)。
相除解得
\[
 \boxed{x=L\,\frac{R_{AB}}{R_{AB}+R_{CD}}
 =5\,\frac{30}{100}=1.50\,\mathrm m}.
\]
即故障点离左端1.5 m、离右端3.5 m。
若存在多点短路或绝缘层漏电，关系不再成立。

\paragraph{习题 16.7}
(1) 局部宏观欧姆定律
\(\boldsymbol j=\boldsymbol E/\rho\)，
所以
\[
 \boxed{\rho=\frac Ej
 =\frac{120}{4\times10^{-12}}
 =3.0\times10^{13}\,\Omega\mathrm m}.
\]
(2) 若进一步\emph{假设}整个全球地表
的平均竖直电流密度均为题给值，
地球半径取 \(R_E=6.37\times10^6\,\mathrm m\)，
则总电流
\[
 I_{\rm global}=j(4\pi R_E^2)
 \approx2.04\times10^3\,\mathrm A.
\]
由电离层—地面电势差
\(U=4.0\times10^5\,\mathrm V\) 求等效电阻
\[
 \boxed{R_{\rm global}=U/I_{\rm global}
 \approx1.96\times10^2\,\Omega}.
\]
只由\emph{局部} \(E,j\)
不能唯一推出全球总电阻，
第二问的数值必须明说全球均匀电流假设。

\paragraph{习题 16.8}
径向电流穿过长 \(L\) 的同轴圆筒壳层，
半径 \(r\) 处有效电流面积为
\(A(r)=2\pi rL\)，
微元电阻
\(dR=\rho\,dr/(2\pi rL)\)。
积分从 \(r_1\) 到 \(r_2\)：
\[
 \boxed{R=\frac{\rho}{2\pi L}
 \ln\frac{r_2}{r_1}
 =\frac{10^{12}\ln4}{2\pi(1000)}
 \approx2.21\times10^8\,\Omega}.
\]
因此电势差100 V时的漏电流
\[
 \boxed{I=U/R\approx4.53\times10^{-7}\,\mathrm A
 =0.453\,\mu\mathrm A}.
\]
忽略轴端效应并假定绝缘层电阻率均匀。

\paragraph{习题 16.9}
同心球之间的径向电流 \(I\) 在任意半径处
流经球面面积 \(4\pi r^2\)，
故局部电场与电势降
\[
 E_r(r)=\rho\,\frac I{4\pi r^2},\qquad
 U=\int_{r_1}^{r_2}E_r\,dr.
\]
由 \(R=U/I\)：
\[
 \boxed{R=\frac{\rho}{4\pi}
 \int_{r_1}^{r_2}\frac{dr}{r^2}
 =\frac{\rho}{4\pi}\left(
 \frac1{r_1}-\frac1{r_2}\right)}.
\]
注意漏电流随 \(1/r^2\) 分布的是\emph{电流密度}，
总电流 \(I\) 对所有同心球面相同。

\paragraph{习题 16.10}
接地点向均匀半无限土壤排入电流 \(I\)，
半径 \(r\) 的半球面面积为
\(2\pi r^2\)，电流密度
\(j=I/(2\pi r^2)\)。
地表电势相对于无穷远为
\[
 \varphi(r)=\int_r^\infty\rho\frac I{2\pi s^2}\,ds
 =\frac{\rho I}{2\pi r}.
\]
所以两脚间跨步电压
\[
 \boxed{|\Delta V|=
 \frac{\rho I}{2\pi}\left(\frac1{1.0}-\frac1{1.3}\right)
 =\frac{2000}{2\pi}\frac{0.3}{1.3}
 \approx73.5\,\mathrm V}.
\]
现实触电风险依人体路径、电缆电压、土壤不均匀性
以及接地几何而变，绝不据此判定现场安全。

\paragraph{习题 16.11}
用原图16.21两个电源的\emph{电动势正端均在下侧}
的阅读方向建模。记共同下结点为
\(V_B=0\)，顶结点为 \(V_A=V\)，
三条从顶到底的电流分别为
\[
 I_1=\frac V{R_1},\quad
 I_m=\frac{V+\mathcal E_1}{r_1},\quad
 I_2=\frac{V+\mathcal E_2}{R_2+r_2}.
\]
结点定律
\(I_1+I_m+I_2=0\)，
代入题给值：
\[
 \frac V2+\frac{V+3}{0.5}
 +\frac{V+6}{5}=0
 \Longrightarrow V=-\frac83\,\mathrm V.
\]
所以外部两电阻上取自顶向下为正的电流：
\[
 \boxed{I_{R_1}=-\frac43\,\mathrm A
 \quad(\text{实际向上})},\quad
 \boxed{I_{R_2}=\frac23\,\mathrm A
 \quad(\text{实际向下})}.
\]
中间电源支路为 \(+2/3\,\mathrm A\) 向下。
\textbf{图示审核风险：}应再对原图16.21
及辅导书复核电池长短板方向；
若极性不同，这一组数值必须重新计算，
不能把假设当作唯一的图示事实。

\paragraph{习题 16.12}
将原图16.22左结点记 \(L\)，
中间结点记 \(M\)，右端上结点记 \(N\)。
上支路 \(L\to N\)
为 \((\mathcal E_1,r_1,R_1)\)，
右端 \(N\to M\) 为 \(R_4\)；
下支路 \(L\to M\)
为 \((\mathcal E_2,r_2,R_2)\)；
横支路 \(L\to M\) 为 \(R_3\)。
上方串联支路总阻为
\(r_1+R_1+R_4=10\,\Omega\)，
下方支路总阻 \(r_2+R_2=20\,\Omega\)。
设教材图16.22两电源的\emph{正端朝右}，
取 \(V_L=0,V_M=V\)，
三个 \(L\to M\) 的电流：
\[
 I_{\rm top}=\frac{\mathcal E_1-V}{10},\quad
 I_{\rm bottom}=\frac{\mathcal E_2-V}{20},\quad
 I_{R_3}=-\frac V{10}.
\]
结点定律
\[
 \frac{3-V}{10}+\frac{1-V}{20}
 -\frac V{10}=0
 \Longrightarrow
 \boxed{V=1.40\,\mathrm V}.
\]
所以
\[
 \boxed{I_{R_1}=I_{R_4}=+0.160\,\mathrm A
 \ (L\to N\to M)},\quad
 \boxed{I_{R_2}=-0.0200\,\mathrm A
 \ (M\to L)},
\]
\[
 \boxed{I_{R_3}=-0.140\,\mathrm A
 \ (M\to L)}.
\]
检验结点：
\(0.160-0.020-0.140=0\)。
本解保留电池图示正端假设，
需用独立高清原图与辅导书核对；
若第二电源极性画反应修改其电动势符号。

\paragraph{习题 16.13}
记电桥上、下四臂电阻按图16.23
为 \(R_1,R_2\)（左边），
\(R_3,R_4\)（右边），
检流计支路电阻 \(R_g\)。
取左端 A 到上节点 B 的电流 \(I_1\)，
左端 A 到下节点 D 的电流 \(I_2\)，
B到D的电流 \(I_g\)，
则右上臂电流 \(I_1-I_g\)，
右下臂为 \(I_2+I_g\)。
三个独立回路关系可写成
\[
 I_1R_1+I_gR_g=I_2R_2,
\]
\[
 (I_1-I_g)R_3=(I_2+I_g)R_4+I_gR_g,
\]
\[
 \mathcal E=I_1R_1+(I_1-I_g)R_3.
\]
在
\(R_1/R_2=R_3/R_4\) 时
解 \(I_g=0\) 一致：
左右分压使
\[
 \frac{V_{AB}}{V_{AD}}=
 \frac{R_1}{R_2}
 =\frac{R_3}{R_4}
 =\frac{V_{BC}}{V_{DC}},
\]
从而 \(V_B=V_D\)，检流计无电压也无电流。
若四臂比不满足，则一般有 \(I_g\ne0\)。

\paragraph{习题 16.14}
依图16.24电流箭头设
\(I_e=I_c+I_b=2.020\,\mathrm{mA}\)，
下端电源正极相对上端为
\(+6.0\,\mathrm V\)，
令顶端电势 \(V_{\rm top}=0\)，
底端 \(V_{\rm bottom}=6.0\,\mathrm V\)。
集电极电流通过
\(R_c=1\,\mathrm{k}\Omega\)，
使
\(V_c=I_cR_c=2.00\,\mathrm V\)。
给
\(U_{ec}=V_e-V_c=1.96\,\mathrm V\)，
故 \(V_e=3.96\,\mathrm V\)；
另给
\(U_{eb}=V_e-V_b=0.20\,\mathrm V\)，
则 \(V_b=3.76\,\mathrm V\)。
因此
\[
 \boxed{R_e=\frac{6.00-3.96}{0.002020}
 \approx1.01\,\mathrm{k}\Omega}.
\]
分压支路电流箭头：
\(I_2=0.400\,\mathrm{mA}\)
由下端流至节点 b，
基极电流 \(I_b=0.020\,\mathrm{mA}\)
从节点 b 向左流入偏置支路，
所以流经 \(R_1\) 的电流
\(I_1=I_2+I_b=0.420\,\mathrm{mA}\)。
故
\[
 \boxed{R_1=\frac{3.76}{0.000420}
 \approx8.95\,\mathrm{k}\Omega},\quad
 \boxed{R_2=\frac{6.00-3.76}{0.000400}
 =5.60\,\mathrm{k}\Omega}.
\]
这是与原图箭头及 \(U_{ec},U_{eb}\)
带符号定义一致的一套结点解；
若图示基极电流朝向在扫描中读取相反，
应重新核对 \(R_1\) 的 \(I_1\) 约定。

\paragraph{习题 16.15}
初始电容电压 \(U_0\)，电荷
\(Q_0=CU_0\)。
理想RC放电：
\[
 Q(t)=Q_0e^{-t/(RC)},\quad
 I(t)=-\dot Q(t)=\frac{U_0}{R}e^{-t/(RC)}.
\]
电阻消耗的热量
\[
 W_R=\int_0^\infty I(t)^2R\,dt
 =\frac{U_0^2}{R}\int_0^\infty
 e^{-2t/(RC)}dt
 =\boxed{\frac12CU_0^2},
\]
恰与电容初储能相同。
数学指数只有在 \(t\to\infty\)
才严格到零，但物理电荷以单个电子为最小粒度，
当残余电量降到 \(e\) 时
\[
 \boxed{t_e=RC\ln\frac{Q_0}{e}
 =RC\ln\frac{CU_0}{e}}.
\]
举例 \(R=1\,\mathrm{k}\Omega\)，
\(C=1\,\mu\mathrm F\)，
\(U_0=100\,\mathrm V\)，
则 \(RC=10^{-3}\,\mathrm s\)，
\(Q_0/e\approx6.24\times10^{14}\)，
故
\(t_e\approx0.0341\,\mathrm s\)。
最后一个电子转移涉及离散性与泄漏，
连续RC模型在如此微小电荷尺度只是近似；
不可据此误说需无穷长时间才能释放实用电能。

\paragraph{习题 16.16}
闪光灯储能电容
\(C=2000\,\mu\mathrm F=0.0020\,\mathrm F\)，
目标电压 \(U_0=4000\,\mathrm V\)。
理想RC恒压充电，初始充电电流最大
\(I_{\max}=U_0/R\)。
若电源允许的最大输出为1 A，
则应取
\(R_{\min}=U_0/I_{\max}=4000\,\Omega\)，
对应\emph{最小充电时间常数}
\[
 \boxed{\tau_{\rm charge,min}
 =R_{\min}C=8.0\,\mathrm s}.
\]
实际上充到精确4000 V的时间在理想RC模型中趋于无穷，
8秒仅是一个时间常数，\emph{不是充满时间}。
闪光灯管内近似电阻 \(R_d=0.5\,\Omega\)，
故最初放电电流最大
\[
 \boxed{I_{\rm discharge,max}=
 \frac{U_0}{R_d}=8.0\times10^3\,\mathrm A},\quad
 \boxed{\tau_{\rm discharge}=R_dC
 =1.0\times10^{-3}\,\mathrm s}.
\]
闪光灯管真实放电电阻随电离过程变化，
此为定电阻近似。

\paragraph{习题 16.17}
电磁铁耗电功率
\(P=UI=(400)(200)=80.0\,\mathrm{kW}\)。
全部转为冷却水获得的热功率时，
\[
 P=\dot m c_p(T_{\rm out}-T_{\rm in}).
\]
允许水温升高
\(\Delta T=80-20=60\,\mathrm K\)，
故所需最小质量流率
\[
 \dot m_{\min}=\frac{80000}{4180(60)}
 \approx0.319\,\mathrm{kg/s}.
\]
取水密度 \(1\,\mathrm{kg/L}\)，
每分钟体积流量
\[
 \boxed{\dot V_{\min}\approx0.319\times60
 =19.1\,\mathrm{L/min}}.
\]
若部分能量通过空气散失，实际所需水量可能下降；
题要求保守的全水冷模型。

\paragraph{习题 16.18}
稳恒电流满足电荷连续方程
\[
 \frac{\partial\rho_f}{\partial t}
 +\boldsymbol\nabla\cdot\boldsymbol j=0.
\]
稳恒时
\(\partial\rho_f/\partial t=0\)，
故
\(\boldsymbol\nabla\cdot\boldsymbol j=0\)。
对于均匀欧姆导体，
\(\boldsymbol j=\sigma\boldsymbol E\)
且导电率 \(\sigma\) 是空间常数，
于是
\[
 0=\boldsymbol\nabla\cdot\boldsymbol j
 =\sigma\,\boldsymbol\nabla\cdot\boldsymbol E
 =\frac{\sigma\rho_f}{\varepsilon_0}.
\]
由 \(\sigma>0\)，
\[
 \boxed{\rho_f=0
 \quad\text{（均匀导体稳恒电流的内部）}}.
\]
界面处 \(\sigma\) 不连续，
\(\nabla\cdot\boldsymbol E\) 仍可支持表面电荷；
而对空间变化的导电率，
\(\nabla\sigma\cdot\boldsymbol E\)
会在连续方程中产生额外项，
不能把证明无条件推广到非均匀材料。
